% Insert as A.46.  This computes the full continuous obstruction of A.43.
\begin{statement}{A.46}{Proved}{The full continuous weight-three residue and its ordinary-support formula}
For every coefficient prime $\ell$, define
\[
 R_\ell=\{(z_q)_{q\ne\ell}\in\prod_{q\ne\ell}\Z_\ell:
                         z_q\longrightarrow0\},\qquad
 (d_3z)_q=(q^2-1)z_q,
 \tag{A.46.1}
\]
where convergence means that each reduction modulo $\ell^N$
has finite support.  The strict localization map induces an
isomorphism
\[
 H^2_{\rm cts}(G_\Q,\Z_\ell(3))
            \xrightarrow{\ \sim\ }R_\ell/d_3R_\ell.
 \tag{A.46.2}
\]
The denominator is the literal image, without taking its closure.

Fix a class-two section represented by $(a,b,h)$ as in A.43.
At $q\ne\ell$, choose arithmetic Frobenius $\Phi_q$ and a
tame $\ell$-inertia generator $T_q$, with
$\Phi_qT_q\Phi_q^{-1}=T_q^q$.  In the fixed Tate coordinates
write
\[
 \begin{array}{lll}
 m_q=a(T_q),&n_q=b(T_q),&t_q=h(T_q),\\
 \alpha_q=a(\Phi_q),&\beta_q=b(\Phi_q),&\gamma_q=h(\Phi_q),
 \end{array}
 \qquad
 \epsilon_q=\frac{q-1}{2},\quad
 e_q=\frac{q^2-1}{24}.
 \tag{A.46.3}
\]
Define the integer-valued function
\[
 \mathcal B_q(z)=q^{-1}\sum_{j=1}^{q-1}\binom{jz}{2}
       =\frac{q-1}{12}\bigl((2q-1)z^2-3z\bigr)
                  \quad(z\in\Z_\ell).
 \tag{A.46.4}
\]
Its first expression proves integrality, because $q$ is a
$\Z_\ell$-unit.  Suppressing the subscript $q$ on the six
variables in the following formulas, set
\[
 \begin{split}
 Z_{D,q}={}&m\gamma-q(\alpha+\epsilon_q)t
       -n\binom\alpha2+\frac{\beta}{q}\binom{qm}{2}
       +\epsilon_q mt+n\mathcal B_q(m),\\
 Z_{F,q}={}&n(\gamma-\alpha\beta)
       +q(\beta+\epsilon_q)(mn-t)+m\binom\beta2-e_q n\\
       &\hspace{1em}-\frac{\alpha}{q}\binom{qn}{2}
       -\epsilon_q(mn^2-nt)-m\mathcal B_q(n).
 \end{split}
 \tag{A.46.5}
\]
These are integral families in $R_\ell$.  Under (A.46.2),
the two actual continuous weight-three obstructions are
\[
                   o_D(s)\longmapsto[Z_D],\qquad
                   o_F(s)\longmapsto[Z_F].
 \tag{A.46.6}
\]
This specifies their entire continuous classes, including a
possible derived-limit contribution.

On the ordinary locus with $v_q(x_q),v_q(1-x_q)\ge0$ at every
finite prime, the same classes have a smaller representative.
Here $m_q=v_q(x_q)$ and $n_q=v_q(1-x_q)$, and at most one
of these integers is positive.  For $m_q>0$ there is a unique
$v_q\in\Z_\ell$ with $\beta_q=(q-1)v_q$; for $n_q>0$ there
is a unique $u_q\in\Z_\ell$ with $\alpha_q=(q-1)u_q$.
Define
\[
 P_{D,q}=
 \begin{cases}
 m_q(\gamma_q-q\alpha_qv_q),&m_q>0,\\
 0,&m_q=0,
 \end{cases}
 \qquad
 P_{F,q}=
 \begin{cases}
 n_q(\gamma_q+\beta_qu_q-e_q),&n_q>0,\\
 0,&n_q=0.
 \end{cases}
 \tag{A.46.7}
\]
Then $P_D,P_F\in R_\ell$ and
\[
                      [Z_D]=[P_D],\qquad [Z_F]=[P_F]
                            \quad\text{in }R_\ell/d_3R_\ell.
 \tag{A.46.8}
\]
In particular the zero locus of the global weight-three
obstruction is exactly the locus where, for every $\ell$,
the two coordinatewise quotients
\[
          \left(\frac{P_{D,q}}{q^2-1}\right)_{q\ne\ell},
          \qquad
          \left(\frac{P_{F,q}}{q^2-1}\right)_{q\ne\ell}
 \tag{A.46.9}
\]
belong to $R_\ell$.  Membership includes both their integrality
and their convergence to zero.
\end{statement}
\begin{proof}
\emph{The unramified term and the surjective residue map.}
Write $G_S=G_{\Q,\{\ell,\infty\}}$ and $M=\Z_\ell(3)$.
For odd $\ell$, the arithmetic \'etale comparison and
\cite[VI, Theorem 8.2]{Weibel2013} identify
$H^2_{\rm cts}(G_S,M)$ with the $\ell$-primary part of
$K_4(\Z)$.  This is zero by
\cite[Theorem 0.1]{BIR_Rognes2000}.

At two, put $S'=\operatorname{Spec}\Z[i,1/2]$.
Its Picard group is zero, and its Brauer group is zero:
the ring is a principal ideal domain, and the only omitted
finite place is the single place over two, while the number
field has no real places; the global invariant sum forces
that one invariant to vanish.  The Kummer sequence gives
$H^2(S',\F_2)=0$.  The successive coefficient extensions for
$M/2^N$ and the arithmetic \'etale comparison give
$H^2(G_{\Q(i),S'},M/2^N)=0$ for all $N$.
The finite $H^1$ groups have zero $\varprojlim^1$, so the
continuous integral $H^2$ over $\Q(i)$ is zero as well.
The corresponding unramified group has index two in $G_S$;
restriction followed by corestriction therefore shows that
$H^2_{\rm cts}(G_S,M)$ is killed by two.  Reduction injects
it into $H^2(G_S,\F_2)$.

The latter group is $\Z/2$ and restriction to $\R$ is an
isomorphism.  Indeed the Kummer sequence over $\Z[1/2]$
identifies it with the two-torsion in its Brauer group;
the allowed local invariants at $2$ and $\infty$ are either
both zero or both $1/2$.  These are the same Picard/Brauer
computations used in A.27, with the global invariant sequence
of \cite[II, Proposition 2.1]{BIR_MilneADT}.
On the other hand complex conjugation acts by $-1$ on
$\Z_2(3)$, and
\[
                    H^2(G_\R,\Z_2(3))=0
\]
because the module has zero invariants.  Commutativity of
restriction and reduction now makes the injective reduction
of $H^2_{\rm cts}(G_S,M)$ zero.  We have proved, at every
coefficient prime,
\[
                         H^2_{\rm cts}(G_S,M)=0.
 \tag{A.46.10}
\]

The coefficient-compatible localization triangle of A.27
now has residue coefficient $M(-1)=\Z_\ell(2)$, on which
Frobenius at $q$ acts by $q^2$.  Its derived inverse-limit
residue complex is exactly
$[R_\ell\xrightarrow{d_3}R_\ell]$ in the same degrees as
there.  The long exact sequence and (A.46.10) give
\[
 0\longrightarrow H^2_{\rm cts}(G_\Q,M)
 \longrightarrow R_\ell/d_3R_\ell
 \longrightarrow H^3_{\rm cts}(G_S,M)
 \longrightarrow H^3_{\rm cts}(G_\Q,M).
 \tag{A.46.11}
\]
For odd $\ell$, the third group is zero.  To check the case
$\ell=2$ without discarding real cohomology, use
\cite[I, Theorem 4.10(c) and Corollary 4.15]{BIR_MilneADT}
and the inverse-limit comparison of
\cite[(1.6) and Theorem 2.2]{BIR_Jannsen1988}.
Finite-coefficient $H^3$ is real cohomology, and its
reduction maps are the maps on
\[
 H^3(G_\R,M/2^N)=(M/2^N)/2(M/2^N)\simeq\Z/2.
\]
These maps are isomorphisms.  The finite $H^2$ groups have
zero $\varprojlim^1$, so restriction identifies
$H^3_{\rm cts}(G_S,M)$ with
$H^3(G_\R,M)=\Z/2$.  Its composite through
$H^3_{\rm cts}(G_\Q,M)$ is this isomorphism, and hence
the last map in (A.46.11) is injective.  The residue map
is therefore surjective also at two.  This proves (A.46.2)
with the exact denominator in (A.46.1).

\emph{The strict family before taking local cohomology.}
Use the actual two-cocycles $W_D,W_F$ of (A.43.1).
They descend on each local group to the tame quotient used
in (A.43.13).  To see this at the cochain level, the
class-two section is a cocycle with values in the pro-$\ell$
Heisenberg group.  The discarded prime-to-$\ell$ inertia
acts trivially and its cocycle restriction is a homomorphism
to that pro-$\ell$ group, hence is zero.  The cocycle
therefore factors through its quotient; so do its three
coordinates and the cyclotomic cochains $\varepsilon,\eta$.
For arbitrary local cochains the strict averaging and
normalization maps of A.27 give the same passage.

The values forced by the class-two cocycle identity are
\[
 \begin{array}{lll}
 a(\Phi T)=\alpha+qm,&
 b(\Phi T)=\beta+qn,&
 h(\Phi T)=\gamma+q^2t+\alpha qn,\\
 a(\Phi^{-1})=-\alpha/q,&
 b(\Phi^{-1})=-\beta/q,&
 h(\Phi^{-1})=(\alpha\beta-\gamma)/q^2,\\
 a(T^j)=jm,& b(T^j)=jn,&
 h(T^j)=jt+\binom j2mn.
 \end{array}
 \tag{A.46.12}
\]
The same identity applied to $\Phi T\Phi^{-1}=T^q$ gives
\[
 (q-1)t+\alpha n-\beta m-\frac{q-1}{2}mn=0.
 \tag{A.46.13}
\]
Substitute (A.46.12) into the strict residue expression
(A.43.13) for $W_D$ and $W_F$.  The sums on inertia are
\[
 \begin{split}
 \sum_{j=1}^{q-1}W_D(T^j,T)
       &=-mt\binom q2-n\sum_{j=1}^{q-1}\binom{jm}{2},\\
 \sum_{j=1}^{q-1}W_F(T^j,T)
       &=(mn^2-nt)\binom q2
                         +m\sum_{j=1}^{q-1}\binom{jn}{2}.
 \end{split}
\]
The two middle differences in (A.43.13) are, respectively,
\[
 qm\gamma+\beta\binom{qm}{2},\qquad
 qn\gamma-\alpha\left(qn\beta+\binom{qn}{2}\right).
\]
Combining these with the first terms gives exactly (A.46.5).
Thus (A.46.5) computes the strict residue representative of
the global continuous cocycle itself.  The symbolic program
\texttt{weight3\_residue\_certificate.py} checks these
identities and the reductions below.

For each coefficient level the global continuous cochains
factor through a finite Galois quotient, after increasing
the level used for the integer-valued binomial functions
when needed.  Outside its finite ramification set they
are literally unramified, and their strict residues are
zero.  Hence the family (A.46.5) has finite support modulo
every $\ell^N$, and belongs to $R_\ell$.
Coefficient compatibility of the strict residue maps proves
(A.46.6) in the derived inverse-limit quotient itself.
This step does not replace that quotient by the product
of the local cohomology quotients.

\emph{Reduction on ordinary nonnegative valuation patterns.}
Suppose $m>0$ and $n=0$.  Then
$1-x_q\in1+q\Z_q$, whose $\ell$-adic Kummer class is zero
because this principal-unit group is pro-$q$ and $q\ne\ell$.
Thus the local cocycle $b$ is the coboundary
$(\chi-1)v$ for a unique $v\in\Z_\ell$.
In particular $\beta=(q-1)v$, and (A.46.13) gives $t=mv$.
Substitution into (A.46.5) yields the exact integral identities
\[
 \begin{split}
 Z_D&=m(\gamma-q\alpha v)
                      +(q^2-1)\binom m2v,\\
 Z_F&=-(q^2-1)m\binom{v+1}{2}.
 \end{split}
 \tag{A.46.14}
\]
For $n>0,m=0$, the class of $x_q\in1+q\Z_q$ is zero,
so $a=(\chi-1)u$, $\alpha=(q-1)u$ and $t=-nu$.
The corresponding identities are
\[
 \begin{split}
 Z_D&=(q^2-1)n\binom{u+1}{2},\\
 Z_F&=n(\gamma+\beta u-e_q)
                      -(q^2-1)\binom n2u.
 \end{split}
 \tag{A.46.15}
\]
When $m=n=0$, (A.46.13) gives $t=0$, and both residues
in (A.46.5) vanish exactly.

These formulas also specify the coordinate changes that
produce the linear terms in (A.46.7).  At $m>0$, conjugation
by $(0,v,0)$ in $H(\Z_\ell)$ sends $(a,b,h)$ to
$(a,0,h-\chi av)$; its central Frobenius value is
$\gamma-q\alpha v$.  At $n>0$, conjugation by $(u,0,0)$
sends it to $(0,b,h+ub)$; subtracting the actual tangential
cocycle at $1$, whose central coordinate is $\eta$, gives
$\gamma+\beta u-e_q$.  The latter tangent has this coordinate
because the path cocycle $f_g$ has coefficient $\eta(g)$
in (A.43.9).  Thus these linear terms come from explicit
conjugations in the actual arithmetic class-two group.

The discarded multiples of $q^2-1$ in (A.46.14)--(A.46.15)
are $d_3$ of families in $R_\ell$.  Indeed $m_q,n_q\to0$
by continuity of the two Kummer cocycles; $u_q,v_q$ are
integral; the functions $\binom{u_q+1}{2}$ and
$\binom{v_q+1}{2}$ are integral; and
\[
 v_\ell\!\binom{m_q}{2}\ge v_\ell(m_q)-v_\ell(2),
 \qquad
 v_\ell\!\binom{n_q}{2}\ge v_\ell(n_q)-v_\ell(2).
\]
Thus every correcting family tends to zero, including at
two.  This proves (A.46.8) with a literal $d_3R_\ell$
correction, and not merely after passing to each local
quotient.  Since $q^2-1$ is nonzero, division in $\Q_\ell$
gives a unique candidate primitive for each coordinate;
it lies in $R_\ell$ exactly as specified by (A.46.9).
The full profinite obstruction is zero exactly when all
coefficient obstructions are zero, by the product identity
for continuous cochains in A.38.  This completes the proof.
\end{proof}

For the A.44 construction, $m_{q_j}=n_j$ and $n_{r_j}=n_j$
at its new support primes, while the two fixed valuations
are $m_3=n_2=1$; all other corresponding valuations are
zero.  Apply (A.46.7) to its unique candidate $s_*$ from
A.43.  The remaining infinite-support tests are exactly
\[
 \left(\frac{n_j(\gamma_{q_j}-q_j\alpha_{q_j}v_{q_j})}
                    {q_j^2-1}\right)_j,
 \qquad
 \left(\frac{n_j(\gamma_{r_j}+\beta_{r_j}u_{r_j}-e_{r_j})}
                    {r_j^2-1}\right)_j
 \tag{A.46.16}
\]
being integral and tending to zero for every coefficient
prime, omitting its own residue place.  The fixed coordinates
at $2$ and $3$ are integral: the only nonunit denominators
there are the two components killed by the normalization of
A.40; at the other allowed coefficient primes the denominators
are units.  They do not affect convergence of the families.
No claim that (A.46.16) vanishes or has the required divisibility
is made.  It is the explicit scalar form of the unresolved
continuous weight-three calculation for these constructed ghosts.
